% §3 Preliminaries — draft §2, statements verbatim from blueprint part 02.
%
% Contents and departures:
%   - notation lead, Defs CM/BF, the toolbox, the "which vocabulary does the work"
%     paragraph (adapted refs), the vanishing lemma with the blueprint's proof;
%   - prop:laplace-continuity / prop:laplace-uniqueness INCLUDED (additive in the
%     blueprint; the draft used both without stating them — the paper states them so
%     that later proofs can cite a numbered statement instead of a bracket);
%   - the blueprint's [T] refinement nodes (def:levy-exponent, the causal cases,
%     transform tightness) are NOT transcribed — formalisation vocabulary, covered by
%     §1.1's account;
%   - the causal vocabulary (causal functions, X₀, T_r, the core 𝒟) LIFTED here from
%     draft §§9–10 heads (NOTES-exposition Part 4 repair 1), so §9 (signaling) stands
%     on this section and cites nothing in Appendix A.

\section{Preliminaries and notation}\label{sec:preliminaries}

Time runs over $\RR$. We write $X = \Lone$, real, with positive cone $X_+$ and
$\int f := \int_{\RR} f(t)\,dt$; $\Mplus(\RR)$ for the finite positive Borel measures on
$\RR$, with total mass $\|\mu\| = \mu(\RR)$; translations $(\shift_a f)(t) = f(t-a)$ for
$a \in \RR$; dilations on functions $(\dil_\sigma f)(t) = \sigma^{-1} f(t/\sigma)$,
$\sigma > 0$, which preserve mass, and on measures $\dil_\sigma \mu :=$ the pushforward of
$\mu$ under $t \mapsto \sigma t$; and convolution
$(\mu * f)(t) = \int f(t - r)\,\mu(dr)$ for $\mu \in \Mplus(\RR)$, $f \in X$. Then
$\|\mu * f\|_1 \le \|\mu\|\,\|f\|_1$, and dilation is compatible with convolution:
\[
  \dil_\sigma(\mu * f) = (\dil_\sigma \mu) * (\dil_\sigma f).
\]

For $\mu \in \Mplus$ supported on $[0,\infty)$ the Laplace transform is
\[
  \hat\mu(s) = \int_{[0,\infty)} e^{-st}\,\mu(dt), \qquad s \ge 0 .
\]
If $\mu \ne 0$ then $\hat\mu(s) > 0$ for all $s \ge 0$.

% shared with blueprint def:completely-monotone
\begin{definition}[completely monotone]\label{def:completely-monotone}
A function $f : (0,\infty) \to \RR$ is \emph{completely monotone} (CM) if $f \in C^\infty$
and $(-1)^n f^{(n)} \ge 0$ for all $n \ge 0$.
\end{definition}

% shared with blueprint def:bernstein-function
\begin{definition}[Bernstein function]\label{def:bernstein-function}
A function $g : (0,\infty) \to [0,\infty)$ is a \emph{Bernstein function} if $g \in C^\infty$
and $g'$ is completely monotone. We write $g \in \BF$, and $g \in \BFz$ if moreover
$g(\zp) = 0$.
\end{definition}

The toolbox that connects the two classes to Laplace transforms is classical, and it is
cited rather than reproven: \sscite{feller2009introduction}, Ch.~XIII, and
\sscite{schilling2012bernstein}, Ch.~1 and~3.

% shared with blueprint prop:bernstein-toolbox
\begin{proposition}[the Bernstein-function toolbox]\label{prop:bernstein-toolbox}
\begin{enumerate}
\item (Bernstein--Widder / Feller criterion.) $f$ is the Laplace transform of a probability
measure on $[0,\infty)$ if and only if $f$ is completely monotone and $f(\zp) = 1$. More
generally, $f$ is completely monotone if and only if $f(s) = \int_{[0,\infty)} e^{-st}\,F(dt)$
for a measure $F$ on $[0,\infty)$, not necessarily finite, and $F$ is then unique.
\item For $g : (0,\infty) \to [0,\infty)$ the following are equivalent: $g \in \BF$; the
composition $h \circ g$ is completely monotone for every completely monotone $h$; and
$e^{-\tau g}$ is completely monotone for every $\tau > 0$.
\item Every $g \in \BF$ is nonnegative, nondecreasing and concave, and has the
Lévy--Khintchine representation
\[
  g(s) = a + b s + \int_0^\infty \bigl(1 - e^{-st}\bigr)\,\nu(dt), \qquad
  a, b \ge 0, \quad \int_0^\infty (1 \wedge t)\,\nu(dt) < \infty,
\]
with $g(\zp) = a$. The triple $(a, b, \nu)$ is unique.
\item $\BF$ is a convex cone, closed under pointwise limits: if $g_n \in \BF$ and
$g_n(s) \to g(s) < \infty$ for every $s > 0$, then $g \in \BF$. Consequently $\BF$ is closed
under mixtures $s \mapsto \int g_u(s)\,m(du)$ with $m \ge 0$, whenever the integral is finite.
\end{enumerate}
\end{proposition}

\emph{Which vocabulary does the work.} Proposition~\ref{prop:bernstein-toolbox} is used as a
bridge and not as a toolkit. The object the arguments below manipulate is the
\emph{representation}, part~(3) --- a drift and a Lévy measure --- and almost every proof in
\S\S\ref{sec:representation}--\ref{sec:signaling} works with it directly: the increments of
Theorem~\ref{thm:increments-bernstein} are produced as representations, not as pointwise
limits of Bernstein functions; the potential kernel of the appendix is constructed as a
subordinator's potential measure, not represented from complete monotonicity; the Mellin
calculus of \S\ref{sec:signaling} mentions neither class. Complete monotonicity and the
derivative-sign definition of $\BF$ are genuinely needed at exactly three places, and it is
worth naming them: Lemma~\ref{lem:selfdecomposable-exponents}'s implications
(1)~$\Rightarrow$~(2)~$\Rightarrow$~(3), where $B(s) = sF'(s)$ has no meaning until the
representation is in hand; Proposition~\ref{prop:pair-regularity}(3) and
Proposition~\ref{prop:thorin-subclass}, where the $\CM$ class is the \emph{subject} rather
than a tool; and Proposition~\ref{prop:bernstein-toolbox} itself, which identifies the two
vocabularies once so that no later proof has to. A reader who wants the shortest path to
Theorem~\ref{thm:main-characterization} or Theorem~\ref{thm:signaling-form} may take~(3) as
the definition of admissibility and never invoke Bernstein--Widder.

The same class carries a second name, after its representation rather than its derivative
signs; several statements below are given in both forms, and
Proposition~\ref{prop:bernstein-toolbox}(3) is the identification.

% shared with blueprint def:levy-exponent
\begin{definition}[Lévy exponent]\label{def:levy-exponent}
A function $g : [0,\infty) \to [0,\infty]$ is a \emph{Lévy exponent} if
\[
  g(s) = b_0\,s + \int_0^\infty \bigl(1 - e^{-st}\bigr)\,\nu(dt),
  \qquad b_0 \ge 0,
\]
for a drift $b_0$ and a positive Borel measure $\nu$ on $(0,\infty)$ --- with no killing term,
so that $g(0) = 0$. We write $g \in \LE$, and call $(b_0, \nu)$ its \emph{triple}.
\end{definition}

By Proposition~\ref{prop:bernstein-toolbox}(3), $\LE = \BFz$ as classes of functions, with
the triple unique. Both names are kept because the two readings matter separately: the proofs
of this paper argue mostly in the representation, and the machine-checked development
(\S\ref{sec:machine-checked}) argues \emph{only} there --- it never defines complete
monotonicity at all --- so a statement whose conclusion is given in $\LE$ is one the
formalisation can state verbatim.

% shared with blueprint lem:vanishing
\begin{lemma}[vanishing lemma]\label{lem:vanishing}
If $g \in \BFz$ and $g(s_0) = 0$ for some $s_0 > 0$, then $g \equiv 0$.
\end{lemma}

\begin{proof}
$g$ is nonnegative and concave by Proposition~\ref{prop:bernstein-toolbox}(3), with
$g(\zp) = g(s_0) = 0$. A nonnegative concave function vanishing at both ends of $[0, s_0]$
vanishes on it, so $g \equiv 0$ on $(0, s_0]$ and hence $g' \equiv 0$ on $(0, s_0)$. Now $g'$
is nonincreasing, being the derivative of a concave function, and nonnegative, being
completely monotone; a nonnegative nonincreasing function that vanishes on $(0,s_0)$ vanishes
identically. So $g' \equiv 0$ on $(0,\infty)$, and with $g(\zp) = 0$ this gives $g \equiv 0$.
\end{proof}

Two further classical facts are used repeatedly and are recorded here so that each later
use can cite a statement; both are from \sscite{feller2009introduction}, Vol.~2,
\S~XIII.1, which states the probability and the general-measure forms as separate theorems
(Theorems~2 and~2a, and Theorems~1 and~1a, respectively).

% shared with blueprint prop:laplace-continuity
\begin{proposition}[continuity theorem for Laplace transforms]\label{prop:laplace-continuity}
Let $\mu_n, \mu$ be probability measures on $[0,\infty)$. Then $\mu_n \to \mu$ weakly if and
only if $\hat\mu_n(s) \to \hat\mu(s)$ for every $s > 0$. The same equivalence holds for
arbitrary measures on $[0,\infty)$, provided that in the direction from transforms to measures
the masses $\hat\mu_n(a)$ are bounded for some $a > 0$.
\end{proposition}

% shared with blueprint prop:laplace-uniqueness
\begin{proposition}[uniqueness of the Laplace transform]\label{prop:laplace-uniqueness}
A measure on $[0,\infty)$ is determined by its Laplace transform: if $\hat\mu(s) = \hat\rho(s)$
for all $s$ in some interval $(a,\infty)$, then $\mu = \rho$.
\end{proposition}

Finally, the causal vocabulary, used from \S\ref{sec:signaling} on and in the appendix. A
function $f \in X$ is \emph{causal} if it vanishes a.e.\ on $(-\infty, 0)$; its transform
$\hat f(s) = \int_0^\infty e^{-st} f(t)\,dt$ is then defined for $s > 0$ with
$|\hat f(s)| \le \|f\|_1$. We write
\[
  \Xz := L^1(\RR_+),
\]
identified with the causal elements of $X$; by (A3) every $\Phi_{x,y}$ restricts to $\Xz$.
On $\Xz$, $(T_r)_{r \ge 0}$ denotes the \emph{delay semigroup},
$(T_r f)(t) = f(t - r)\,\mathbf{1}_{t \ge r}$, a strongly continuous semigroup of positive
isometries, and
\[
  \core \;:=\; \bigl\{\, f \in \Xz :\ f \text{ absolutely continuous},\ f' \in \Xz,\
  f(0) = 0 \,\bigr\}
\]
is the operator core used there.
